% SPDX-License-Identifier: CC-BY-SA-4.0
% Part: Introduction
% Section: Words and languages
% Exercise: Revision exercise

\begingroup

\begin{exercice}[Révisions -- Preuves sur les langages]\label{exo:introduction/languages/revision}
  
  Soit $\Sigma = \{a,b,c\}$ et le mot $x = bacca$.

  \begin{question}
  \item Donner la valeur de $|x|$, $|x|_a$ et $|x|_c$.
  \item Donner l’ensemble $\mathit{Pref}(x)$ et l’ensemble $\mathit{Suff}(x)$.
  \item Donner tous les facteurs de $x$ de longueur $2$.
  \end{question}

  Soient les langages $L_1 = \{a,ba,ca\}$ et $L_2 = \{\varepsilon,b,ac\}$ sur l’alphabet $\Sigma$.

  \begin{question}
  \item Calculer $L_1 \cup L_2,$ $L_1 \cap L_2$, $L_1 \cdot L_2$, $L_2 \cdot L_1$ et $L_1^2$.
  \item Donner $\Sigma^\star \setminus L_1$ en le décrivant en français.
  \end{question}

  Le but de la suite des questions est de démontrer que pour tout langage $L$, $(L^\star)^\star = L^\star$.
  Soient $\Sigma$ un alphabet, et $L \in \mathscr{P}(\Sigma^\star)$ un langage sur $\Sigma$.  

  \begin{question}
  \item Montrer, par double inclusion, que $L^\star \times L^\star = L^\star$. 
  \item Montrer, par récurrence simple sur $n$, que $\forall n \ge 1, (L^\star)^n = L^\star$.
  \item En déduire que $(L^\star)^+ = L^\star$.
  \item En déduire que $(L^\star)^\star = L^\star$.
  \end{question}

\end{exercice}

\begin{correction}

  \begin{question}

  \item $|x| = 5$, \quad $|x|_a = 2$, \quad $|x|_c = 2$ \quad (pour $x = bacca$).

  \item $\mathit{Pref}(x) = \{\varepsilon,\ b,\ ba,\ bac,\ bacc,\ bacca\}$ \\
        $\mathit{Suff}(x) = \{\varepsilon,\ a,\ ca,\ cca,\ acca,\ bacca\}$

  \item Les facteurs de longueur $2$ de $x$ (sous-mots contigus) sont :\\
        $\{\,ba,\ ac,\ cc,\ ca\,\}$

  \item 
    $$
    \begin{aligned}
      L_1 \cup L_2 \;&=\; \{a,ba,ca,\ \varepsilon,b,ac\} \\
      L_1 \cap L_2 \;&=\; \emptyset \\
      L_1 \cdot L_2 \;&=\; \{a,ab,aac,\ ba,bab,baac,\ ca,cab,caac\} \\
      L_2 \cdot L_1 \;&=\; \{a,ba,ca,\ bba,bca,\ aca,acba,acca\} \\
      L_1^2 \;&=\; \{aa,aba,aca,\ baa,baba,baca,\ caa,caba,caca\}
    \end{aligned}
    $$

  \item $\Sigma^\star \setminus L_1$ est \textit{l’ensemble de tous les mots sur $\{a,b,c\}$ qui ne sont pas $a$, $ba$ ou $ca$}, c’est-à-dire
        tous les mots de $\Sigma^\star$ sauf ces trois-là.

  \item 
    \begin{description}
    \item[Preuve de $L^\star \subseteq L^\star \times L^\star$ :] Soit $u \in L^\star$. On a $u \in L^\star$ et $\varepsilon \in L^\star$, donc
      $u = u \cdot \varepsilon \in L^\star \times L^\star$
    \item[Preuve de $L^\star \times L^\star \subseteq L^\star$ :] Soit $u \in L^\star \times L^\star$.\\
      Par définition de la concaténation de langages, il existe $v \in L^\star$ et $w \in L^\star$ tels que $u=v\cdot w$.\\
      Par définition de la fermeture de Kleene, il existe $n$ et $m$ tels que $v \in L^n$ et $w \in L^m$.\\
      On a donc $u=v\cdot w \in L^n \times L^m = L^{n+m}$, donc $u\in L^\star$.
    \item[Conclusion :] On a $L^\star \subseteq L^\star \times L^\star$ et $L^\star \times L^\star \subseteq L^\star$,
      donc par double inclusion, $L^\star \times L^\star = L^\star$. 
    \end{description}
  \item Posons le prédicat $P(n) \eqdef (L^\star)^n = L^\star$.
    \begin{description}
    \item[Initialisation :] Pour $n = 1$, on a $(L^\star)^1 = L^\star$, donc $P(1)$.  
    \item[Hérédité :] Soit $n \ge 1$ tel que $P(n)$. Montrons $P(n+1)$. On a:
      $$\begin{array}{rcll}
      (L^\star)^{n+1} &=& (L^\star)^{n} \cdot (L^\star) & \text{Par définition de $(L^\star)^{n+1}$}\\
      &=& L^\star \cdot L^\star & \text{Par $P(n)$}\\
      &=& L^\star & \text{D'après la question 1}\\
      \end{array}$$
    \item[Conclusion :] Par le principe de récurrence, nous avons montré $P(n)$ pour tout $n\ge 1$.
    \end{description}
  \item On a :
      $$\begin{array}{rcll}
      (L^\star)^+ &=& \bigcup_{n\ge 1} (L^\star)^{n} & \text{Par définition de $(L^\star)^+$}\\
      &=& \bigcup_{n\ge 1} L^\star & \text{D'après la question 2}\\
      &=& L^\star & \text{Par idempotence de l'union}\\
      \end{array}$$
  \item On a $(L^\star)^\star = \{\varepsilon\} \cup (L^\star)^+ = \{\varepsilon\} \cup L^\star = L^\star$ car $\varepsilon\in L^\star$.
  \end{question}

\end{correction}

\endgroup
\endinput
